%%%%%%%%%%%%%%%%%%%%%%%%%%%%%%%%%
%
%       EXERCÍCIO
%
%%%%%%%%%%%%%%%%%%%%%%%%%%%%%%%%%

\definevalorquestao

\begin{exercicioBanco}[\valorquestao]
Para que o rendimento de uma síntese orgânica seja otimizado em um reator descontínuo, o tempo de reação deve estar entre $45$ e $48$ minutos. Devido a oscilações na pureza dos reagentes e na transferência de calor, o tempo total de reação pode ser modelado por uma variável aleatória com distribuição normal de média $47{,}4$ minutos e desvio padrão de $1{,}2$ minuto. O setor de controle de qualidade seleciona uma amostra aleatória de $9$ bateladas e só aprova a produção do dia se o tempo médio de reação da amostra estiver no intervalo aceito de $45$ a $48$ minutos. Supondo que as $9$ bateladas sejam independentes, calcule a probabilidade de aprovação do lote.
\end{exercicioBanco}

